\section{Threats to validity}
\label{sec:threats}

This section collects what would have to be true for the conclusions above to
be wrong. They are gathered here rather than scattered through the report so
that a reader can weigh them together, and each names a specific confound and
the measurement that would settle it, rather than hedging.

\paragraph{The split is temporal, not random.}
The recordings held back are the last seven of a single session, because that is
what the trainer withholds. Any property that drifted over the session sits
entirely in the held-back set: the lighting, where the garment was placed
between takes, wear and residue on the tactile gels. The generalisation gap
reported in Section~\ref{sec:results} is therefore a gap in time as well as a
gap in familiarity, and the two cannot be separated from these runs. What would
settle it is one further retraining per family on a randomly chosen seven, with
everything else held fixed; if the gap persists it is overfitting, and if it
collapses it was drift. This is the single largest open question in the report,
and it bears on its largest effect.

\paragraph{The crop changes two things at once.}
Cropping removes the outer rows of the tactile image and then resizes what is
left back to the original shape, which stretches it vertically by a quarter.
Every statement in this report about cropping is therefore a statement about
rim removal \emph{and} anisotropic stretching together. An arm that crops
without resizing would separate them. Until that exists, a null result is a
null result for this operation, and not evidence that the rim carries no
information.

\paragraph{The families do not read the same cameras.}
One family finetunes a pretrained model with a fixed number of image slots,
fewer than the rig has cameras, so it sees the overhead view and two of the
four tactile sensors while the others see all five. Its row is consequently not
an ablation of the same input as the rows above it, and a difference between
that family and the others confounds architecture with input set.

\paragraph{The held-out sample is small.}
The held-out errors rest on \num{111} scored frames. A difference of a few per
cent across that sample is not distinguishable from sampling noise, which
covers the crop effects reported in Section~\ref{sec:cropping}; a difference of
severalfold is far outside it, which covers the generalisation gaps. Where a
claim falls on the wrong side of that line, Section~\ref{sec:findings} says so
rather than reporting the number as though its size were established.

\paragraph{Attribution measures a difference, not an importance.}
Suppressing an input and measuring how much the planned action changes says how
much that input is being used by this trained policy on these frames. It does
not say the input is necessary for the task, nor that a policy trained without
it would do worse --- a stream can be ignored because a correlated one carries
the same information. Retraining with the stream removed is what answers the
second question, and it costs a training run per stream rather than minutes.

\paragraph{One family's numbers are still pending.}
The matched retraining of the third family had not finished when this was
written, so its rows are dashes. Every cross-family statement here is therefore
a statement about two families, and the third may not follow the same pattern:
it differs in the cameras it reads, in being a finetune of a pretrained model
rather than a model trained from scratch, and in planning a different distance
ahead again.

\paragraph{None of this is an on-robot result.}
Error against recorded commands measures agreement with one demonstrator on
recordings that demonstrator produced. It cannot measure whether a policy
recovers from a state the demonstrator never entered, and a lower error is not
the same as a higher chance of folding the garment. Only trials on the rig
settle that, and they would not replace these numbers --- they would bound what
the numbers are evidence for.
